\section*{\texorpdfstring{\S\CurrentExerciseGroup}{第{\CurrentExerciseGroup}节}}
\phantomsection
\addcontentsline{toc}{section}{第{\CurrentExerciseGroup}节习题}

¶ 1) 设 I 是半开区间 \([a, b]\)（\(\mathbf{R}\) 中），并设 \(F = R^I\) 为 I 到 \(\mathbf{R}\) 的全体映射所构成的空间，赋以逐点收敛拓扑；对每个 \(x \in I\)，记 \(\mathbf{f}(x)\) 为映射 \(t \mapsto |x - t|\)（I 到 \(\mathbf{R}\)），它是 \(F\) 中的一个元素；\(f\) 是 I 到 \(F\) 的连续映射。试证 \(f\) 在 I 的每一点都左可微，但左导数 \(\mathbf{f}_l'\) 是一个（取值于 \(F\) 的）函数，它关于 Lebesgue 测度不可测，尽管它是一列连续函数的逐点极限，且对每个 \(t \in I\)，函数 \(\text{pr}_t \circ \mathbf{f}_l'\) 是可测的数值函数（注意 \(\mathbf{f}_l'\) 在 I 的任何一点都不右连续，并利用 GT, IV, §2 习题 1）。

\begin{enumerate}
\def\labelenumi{\arabic{enumi})}
\setcounter{enumi}{1}
\tightlist
\item
  设 \(\nu\) 是 \(\mathbf{R}\) 上的 Lebesgue 测度，\(g\) 是 \(\mathbf{R}\) 到 \([0,1]\) 的连续映射，其支集含于 \([-1,2[\) 且在 \([0,1]\) 上等于 1；再设 \(\mu\) 为测度 \(g \cdot \nu\)（\(\mathbf{R}\) 上）。§4 习题 8 中定义的集合 \(\mathrm{H}_0 \subset [0,1]\) 不是 \(\mu\) -可测的，但集合 \(\mathrm{H} = \mathrm{H}_0 - 2\) 是 \(\mu\) -可忽略的。
\end{enumerate}

\begin{enumerate}
\def\labelenumi{\alph{enumi})}
\item
  若设 \(f(x) = x - 2\)，则 \(f\) 连续且 \(\varphi_{\mathrm{H}}\) 是 \(\mu\) -可测的，但复合函数 \(\varphi_{\mathrm{H}} \circ f\) 不是 \(\mu\) -可测的。
\item
  若设 \(h(x) = x + 2\)，则 \(h(\mathrm{H})\)（\(\mu\) -可测集 \(\mathrm{H}\) 经连续函数 \(h\) 的像）不是 \(\mu\) -可测的。
\end{enumerate}

\begin{enumerate}
\def\labelenumi{\arabic{enumi})}
\setcounter{enumi}{2}
\tightlist
\item
  设 \(f\) 是 \(X\) 到拓扑空间 \(F\) 的可测映射，\(g\) 是 \(F\) 上的下半连续数值函数；证明 \(g \circ f\) 可测。
\end{enumerate}

¶ 4) 设 \(\mu\) 是 \(X = [0,1[\) 上的 Lebesgue 测度，又设 \((H_n)\) 是 \(X\) 的一个分割，分成无穷一列都具有连续统的势且都不可测的集合，并且使得：对每个并 \(H\)（由有限个集合 \(H_n\) 所成），有 \(\mu_*(H) = 0\)（§4 习题 8）。设 \(\sigma_n\) 是区间 \(]1/(n+1),1/n]\) 到 \(H_n\) 上的一个双射。对每个数 \(y\)（满足 \(0 < y \leqslant 1\)），设 \(n\) 是使 \(1/(n+1) < y \leqslant 1/n\) 的整数。定义 \(f_y\) 为化为一点 \(\sigma_n(y)\) 的集合的特征函数；则对每个 \(x \in X\)，\(f_y(x)\) 当 \(y\) 趋于 0 时趋于 0。试证：不存在任何紧集 \(K \subset X\)（测度 \textgreater0）使得 \(f_y\) 在 \(K\) 上一致趋于 0。

\begin{enumerate}
\def\labelenumi{\arabic{enumi})}
\setcounter{enumi}{4}
\tightlist
\item
  设 \((f_{mn})\) 是 \(X\) 到可度量化空间 \(F\) 的可测映射的二重序列。假设对每个 \(m\)，序列 \((f_{mn})_{n\geqslant 1}\) 局部几乎处处收敛于函数 \(g_{m}\)，且序列 \((g_{m})\) 局部几乎处处收敛于函数 \(h\)。试证：对每个紧子集 \(K\)（\(X\) 的），存在两个严格递增的整数序列 \((m_k),(n_k)\)（整数 \(>0\)），使得函数序列 \(f_{m_k,n_k}\) 几乎处处收敛于 \(h\)（在 \(K\) 中）。（注意，对每个 \(\varepsilon >0\)，存在紧集 \(K_1\subset K\) 使得 \(\mu (K - K_1)\leqslant \varepsilon\)，并且在 \(K_{1}\) 中序列 \((g_{m})\) 以及每个序列 \((f_{mn})_{n\geqslant 1}\) 都一致收敛。）
\end{enumerate}

¶ 6) 对每个整数 \(n \geqslant 1\)，设 \(f_n(x) = [2^n x] - 2[2^{n-1}x]\)。在 \(\mathbf{R}\) 到 \(\{0,1\}\) 的映射所成的紧空间（赋以逐点收敛拓扑）中，设 \(f\) 是序列 \((f_n)\) 的一个聚点。试证：对每个二进数 \(r\)，有 \(f(r+x) = f(x)\)（对所有 \(x \in \mathbf{R}\)），且有 \(f(r-x) = 1 - f(x)\)（对每个异于二进数的 \(x \in \mathbf{R}\)）。由此推出，对 Lebesgue 测度 \(\mu\) 而言 \(f\) 不可测。（用反证法；设 \(A\) 为 \(x \in [0,1]\) 中使 \(f(x) = 1\) 者组成的集，并假设 \(A\) 可测且 \(\mu(A) = \alpha > 0\)；证明存在一个集 \(I\)，它是含于 \(]0,1[\) 内的开区间的有限并，使得 \(\mu(I \cap \mathbb{C}A) \leqslant \alpha/4\) 且 \(\mu(A \cap \mathbb{C}I) \leqslant \alpha/4\)；考察 \(I\) 中形如 \(]k/2^n, (k+1)/2^n[\) 的区间，证明这将与关系式 \(f(r-x) = 1 - f(x)\)（\(r\) 为二进数、\(x\) 为非二进数）矛盾。）

\begin{enumerate}
\def\labelenumi{\arabic{enumi})}
\setcounter{enumi}{6}
\tightlist
\item
  设 \(\mathbf{X}\) 为区间 \([0,1]\)（\(\mathbf{R}\) 中），并设 \(\mathbf{F}\) 为具有标准正交基 \((\mathbf{e}_t)_{0\leqslant t\leqslant 1}\)（与 \(\mathbf{X}\) 等势）的 Hilbert 空间。
\end{enumerate}

\begin{enumerate}
\def\labelenumi{\alph{enumi})}
\item
  证明映射 \(\mathbf{f}\)（\(X\) 到 \(F\)，满足 \(\mathbf{f}(t) = \mathbf{e}_t\)，\(0 \leqslant t \leqslant 1\)）关于 Lebesgue 测度不可测，但 \(\mathbf{f}\) 下 \(F\) 中每个闭球的逆像均可测，且对每个连续线性形式 \(\mathbf{a}'\)（\(F\) 上），\(\langle \mathbf{f}, \mathbf{a}' \rangle\) 都可忽略。
\item
  设 H 是 X 中一个不可测集（§4 习题 8）；证明若 \(\mathbf{g} = \mathbf{f}\varphi_{\mathrm{H}}\)，则函数 \(\langle \mathbf{g},\mathbf{a}'\rangle\) 对每个连续线性形式 \(\mathbf{a}'\)（F 上）均可忽略，但数值函数 \(|\mathbf{g}|\) 不可测。
\end{enumerate}

\begin{enumerate}
\def\labelenumi{\arabic{enumi})}
\setcounter{enumi}{7}
\item
  设 \(\mathbf{F}\) 是可度量化局部凸空间，\(\mathbf{f}\) 是 \(\mathbf{X}\) 到 \(\mathbf{F}\) 的映射，满足第 5 号命题 10 推论 1 中的条件 \(a)\) 与 \(b)\)；证明 \(\mathbf{f}\) 可测。
\item
  设 \(\mu\) 是 \(X = [0,1]\) 上的 Lebesgue 测度；用 \(F\) 表示以 \(\mathbf{R}\) 为基域、由 \(\mu\) -可测的有限数值函数（\(X\) 上）所组成的向量空间，赋以逐点收敛拓扑，这使它成为 Hausdorff 局部凸空间。
\end{enumerate}

\begin{enumerate}
\def\labelenumi{\alph{enumi})}
\item
  证明 \(\mathbf{F}\) 中存在可数稠密子集（在 Banach 空间 \(\mathcal{C}(\mathbf{X};\mathbf{R})\)（由 \(\mathbf{X}\) 上的连续数值函数组成）中取一个稠密序列）。
\item
  对每个 \(x \in \mathbf{X}\)，设 \(\mathbf{f}(x)\) 是对偶 \(\mathrm{F}'\)（即 \(\mathrm{F}\) 的对偶）中由 \(\langle z, \mathbf{f}(x) \rangle = z(x)\)（对一切 \(z \in \mathrm{F}\)）定义的元素。证明当 \(\mathrm{F}'\) 赋以弱拓扑 \(\sigma(\mathrm{F}', \mathrm{F})\) 时，\(\mathbf{f}\) 不是 \(\mu\) -可测的，但 \(\langle a, \mathbf{f} \rangle\) 是 \(\mu\) -可测的（对每个 \(a \in \mathrm{F}\)）。
\end{enumerate}

¶ 10) 设 F 是 Banach 空间，f 是 X 到 F 的可测映射，且 \(x \in X\) 中使 \(\mathbf{f}(x) \neq 0\) 者组成的集 A 是可数个可积集的并。证明存在一列具有紧支集、取值于 F 的连续函数 \((\mathbf{f}_n)\)，使得序列 \((\mathbf{f}_n(x))\) 在 X 中几乎处处收敛于 \(\mathbf{f}(x)\)（一方面注意，A 是一个可忽略集 N 与一列两两不相交的紧集 \(K_n\) 的并，且 f 在每个 \(K_n\) 上的限制连续；另一方面注意，存在包含 A 的递减开集序列 \((\mathrm{U}_n)\)，使得当 n 趋于无穷时 \(\mu(\mathrm{U}_n \cap \mathbb{C}\mathrm{A})\) 趋于 0）。

\begin{enumerate}
\def\labelenumi{\arabic{enumi})}
\setcounter{enumi}{10}
\tightlist
\item
  设 \(\mathbf{f}\) 是 \(X\) 到 Banach 空间 \(F\) 的可测映射。对每个有理整数 \(n\)（正或负），设 \(A_{n}\) 为 \(x \in X\) 中使 \(2^{n} \geqslant |\mathbf{f}(x)| > 2^{n-1}\) 者组成的集。为使 \(\mathbf{f}\) 可积，必要且充分的是通项为 \(2^{n}\mu(A_{n})\)（\(n \in \mathbf{Z}\)）的级数收敛。
\end{enumerate}

¶ 12) 设 \(\mu\) 是 X 上的一个测度 \(\geqslant 0\)。设 \((\mathbf{f}_n)\) 是 X 上的一列可积函数，在 X 上逐点收敛于函数 \(\mathbf{f}\)。

\begin{enumerate}
\def\labelenumi{\alph{enumi})}
\tightlist
\item
  证明若 \(\mathbf{f}\) 可积且
\end{enumerate}

\[
\int \mathbf {f} d \mu = \lim _ {n \rightarrow \infty} \int \mathbf {f} _ {n} d \mu ,
\]

则对每个 \(\varepsilon > 0\)，存在一个可积集 \(A\)、一个可积函数 \(g \geqslant 0\) 和一个整数 \(n_0\)，使得对每个 \(n \geqslant n_0\),

\[
\left| \int \mathbf {f} _ {n} \varphi_ {\mathbb {C} A} d \mu \right| \leqslant \varepsilon
\]

而且 \(|\mathbf{f}_n(x)| \leqslant g(x)\) 对一切 \(x \in A\) 成立（考察一个可积集 \(B\)，使得

\[
\int | \mathbf {f} | \varphi_ {\mathbb {C} _ {\mathrm{B}}} d \mu \leqslant \varepsilon / 2
\]

并使 \(\mathbf{f}\) 在 B 上有界，然后应用 Egoroff 定理）。

\begin{enumerate}
\def\labelenumi{\alph{enumi})}
\setcounter{enumi}{1}
\item
  假设对每个 \(\varepsilon > 0\)，存在一个可测集 \(A\)、一个可积函数 \(g \geqslant 0\) 与一个整数 \(n_0\)，使得对每个 \(n \geqslant n_0\)，有 \(\int |\mathbf{f}_n| \varphi_{\mathbb{C}A} d\mu \leqslant \varepsilon\)，而且 \(|\mathbf{f}_n(x)| \leqslant g(x)\) 对一切 \(x \in A\) 成立。证明在这些条件下，\(\mathbf{f}\) 可积且 \(\widetilde{\mathbf{f}}_n\) 趋于 \(\widetilde{\mathbf{f}}\)（在空间 \(L^1\) 中）。反之亦然。
\item
  设 \(\mathbf{F} = \mathbf{R}\)；试举例说明，\(a\) ) 的条件并不充分，而 \(b\) ) 的条件并非必要（就 \(\mathbf{f}\) 可积且 \(\int \mathbf{f} d\mu = \lim_{n \to \infty} \int \mathbf{f}_n d\mu\) 成立而言）。
\end{enumerate}

¶ 13) 设 \(\mu\) 是 X 上的正测度。若 A 可测，证明对 X 的每个子集 B，

\[
\mu^ {*} (\mathrm{B}) = \mu^ {*} (\mathrm{B} \cap \mathrm{A}) + \mu^ {*} (\mathrm{B} \cap \mathbb {C} \mathrm{A})
\]

（若 \(\mu^{*}(\mathrm{B}) < +\infty\)，考察一个可积集 \(\mathrm{B}_{1}\)，使得 \(\mathrm{B} \subset \mathrm{B}_{1}\) 且 \(\mu^{*}(\mathrm{B}) = \mu(\mathrm{B}_{1})\)（§4 习题 7 b)））。反之，证明若 A 满足此条件，则 A 可测（参见 §4 习题 6 e)）。

\begin{enumerate}
\def\labelenumi{\arabic{enumi})}
\setcounter{enumi}{13}
\tightlist
\item
  设 \((\mathbf{A}_n)\) 是 \(X\) 的子集序列，使得对每个指标 \(n\)，存在可测集 \(\mathbf{B}_n \supset \mathbf{A}_n\)，且诸 \(\mathbf{B}_n\) 两两不相交。证明
\end{enumerate}

\[
\mu^ {*} \left(\bigcup_ {n} \mathrm{A} _ {n}\right) = \sum_ {n} \mu^ {*} (\mathrm{A} _ {n}) \quad \text { and } \quad \mu_ {*} \left(\bigcup_ {n} \mathrm{A} _ {n}\right) = \sum_ {n} \mu_ {*} (\mathrm{A} _ {n}).
\]

\begin{enumerate}
\def\labelenumi{\arabic{enumi})}
\setcounter{enumi}{14}
\tightlist
\item
  设 \(\mu\) 是 \(X\) 上的正测度。数值函数 \(f\)（有限或非有限，定义在 \(X\) 上）称为拟可积的，如果它可测且 \(\mu^{*}(f) = \mu_{*}(f)\)（§4 习题 5）。于是令
\end{enumerate}

\[
\mu (f) = \mu^ {*} (f) = \mu_ {*} (f);
\]

同样，用 \(\int f d\mu\) 代替 \(\mu(f)\) 来书写。

\begin{enumerate}
\def\labelenumi{\alph{enumi})}
\item
  证明：若 \(f\) 与 \(g\) 拟可积且和式 \(\mu(f) + \mu(g)\) 有定义，则 \(f + g\) 几乎处处有定义且拟可积，并且 \(\mu(f + g) = \mu(f) + \mu(g)\)。
\item
  由 a) 推出：为使 \(f\) 拟可积，必要且充分的是 \(f\) 可测且数 \(\mu^{*}(f^{+}), \mu^{*}(f^{-})\) 中至少一个有限；此时 \(\mu(f) = \mu^{*}(f^{+}) - \mu^{*}(f^{-})\)。
\end{enumerate}

¶ 16) 设 X 是紧空间，\(\mu\) 是 X 上的正测度。定义在 X 上的有界数值函数 \(f\) 称为在 X 中（关于测度 \(\mu\) ）几乎处处连续，如果 X 中使 \(f\) 连续（相对于 X）的点集的余集测度为零。

\begin{enumerate}
\def\labelenumi{\alph{enumi})}
\item
  给出一个几乎处处连续的函数 \(f\) 的例子，使得不存在任何连续函数 \(g\) 几乎处处等于 \(f\)。
\item
  假设 \(\mu\) 的支集恰为 X。证明：为使定义在 X 上的有界数值函数 f 几乎处处等于一个在 X 中几乎处处连续的函数，必要且充分的是，存在 X 的一个子集 H，其余集可忽略，使得 f 到 H 的限制 \(f|H\) 连续（为证条件充分，注意 H 在 X 中稠密，且 \(f|H\) 的在 X 上下半连续的延拓（GT, IV, §6, 命题 4）在 H 的每一点都是（相对于 X 的）连续函数）。由此推出 f 可测。
\item
  由 b) 推出：若此外 X 可度量化，则对每个几乎处处等于一个在 X 中几乎处处连续的函数的有界数值函数 f，存在一列在 X 上连续的函数 \((f_{n})\)，它在 X 的每一点都收敛，且其极限几乎处处等于 f（参见 §4 习题 4 c)）（注意该命题对下半连续函数 f 成立）。
\item
  设 A 是一个无内点的完全集，含于 X 且测度 \textgreater0（参见 §4 习题 4 a)）。证明不存在任何在 X 中几乎处处连续且几乎处处等于上半连续函数 \(\varphi_{\mathrm{A}}\) 的函数。
\end{enumerate}

¶ 17) 设 X 是紧空间，\(\mu\) 是 X 上的正测度。称 X 的由可积集组成的有限分割所成的一个集合 \(\mathcal{P}\)（按关系 \(\langle\varpi\) 粗于 \(\varpi'\rangle\) 定向）为基本的，如果对 X 的一致结构的每个邻域 V，存在 X 的一个分割 \(\varpi = (A_i)\)，它属于 \(\mathcal{P}\)，使得所有 \(A_i\) 都是关于 V 的小集。对每个有限分割 \(\varpi = (A_k)\)（属于 \(\mathcal{P}\)）以及 X 上的每个有界数值函数 \(f\)，令

\[
s _ {\varpi} (f) = \sum_ {k} \inf _ {x \in \mathrm{A} _ {k}} f (x) \cdot \mu (\mathrm{A} _ {k}) \quad \text { and } \quad S _ {\varpi} (f) = \sum_ {k} \sup _ {x \in \mathrm{A} _ {k}} f (x) \cdot \mu (\mathrm{A} _ {k})
\]

（关于 \(f\) 与分割 \(\varpi\) 的“Riemann 和”。）

\begin{enumerate}
\def\labelenumi{\alph{enumi})}
\item
  证明 \(s_{\varpi}(f) \leqslant \mu_{*}(f) \leqslant \mu^{*}(f) \leqslant S_{\varpi}(f)\) 对每个分割 \(\varpi \in \mathcal{P}\) 成立，且 \(s_{\varpi}(f)\) 与 \(S_{\varpi}(f)\) 各自沿有向有序集 \(\mathcal{P}\) 趋于一个极限。
\item
  若 \(\mathcal{P}\) 是 X 的由可积集组成的所有有限分割构成的（基本）集合，证明对每个有界且可积的函数 \(f\)，\(s_{\varpi}(f)\) 与 \(S_{\varpi}(f)\) 趋于 \(\int f d\mu\)（沿 \(\mathcal{P}\)）。
\item
  若 \(f\) 是在 \(X\) 中几乎处处连续的有界函数，则 \(s_{\varpi}(f)\) 与 \(S_{\varpi}(f)\) 趋于 \(\int f d\mu\)（沿每个基本集合 \(\mathcal{P}\)，即 \(X\) 的由可积集组成的有限分割的基本集合）（对每个 \(\varepsilon > 0\)，考察闭集 \(A\)（\(f\) 在其上振幅 \(\geqslant \varepsilon\) 的点集），并且对每个分割 \(\varpi \in \mathcal{P}\)（其各集合均为关于 \(V\) 的小集），分别考察 \(\varpi\) 中与 \(V(A)\) 相交的集合以及与其不相交的集合）。若 \(f\) 在 \(X\) 上有界且下半连续，证明 \(s_{\varpi}(f)\) 趋于 \(\int f d\mu\)（沿每个基本集合 \(\mathcal{P}\)，即 \(X\) 的由可积集组成的有限分割的基本集合）（把 \(f\) 看作连续函数的上包络）。
\item
  称集合 \(A \subset X\)（对 \(\mu\) ）为可求积的，如果其特征函数几乎处处连续，换言之，如果其边界是 \(\mu\) -可忽略的。证明 X 的每一点 \(x_{0}\) 都有一个由可求积开邻域组成的基本邻域系（对 \(x_{0}\) 的每个邻域 V，设 f 是取值于 {[}0,1{]} 的连续函数，在点 \(x_{0}\) 处等于 1 而在 CV 上等于 0；考察使 \(f(x) > \alpha\) 的 x 组成的集合，其中 \(0 < \alpha < 1\)）。由此推出，存在 X 的有限分割的一个基本集合 P，使得每个分割 \(\varpi \in P\) 都由开集与可忽略集组成。
\item
  设 \(\mathcal{P}\) 是 \(X\) 的由可积集组成的有限分割的一个基本集合，使得每个分割 \(\varpi \in \mathcal{P}\) 都由开集与可忽略集组成。对每个有界函数 \(f\)（\(X\) 上），设 \(g\) 是 \(X\) 上 \(\leqslant f\) 的最大下半连续函数（GT, IV, §6, 命题 4）；证明对每个分割 \(\varpi \in \mathcal{P}\)，有 \(s_{\varpi}(f) = s_{\varpi}(g)\)。由此推出，为使 \(s_{\varpi}(f)\) 与 \(S_{\varpi}(f)\) 沿 \(\mathcal{P}\) 趋于一个公共极限，必须 \(f\) 在 \(X\) 中几乎处处连续。
\item
  由 e) 推出一个可忽略函数 f 以及 X 的由可积集组成的有限分割的一个基本集合 P 的例子，使得 \(s_{\varpi}(f)\) 与 \(S_{\varpi}(f)\) 沿 P 不趋于同一极限（取 X 为 R 中的区间 {[}0,1{]}，取 \(\mu\) 为 Lebesgue 测度）。
\end{enumerate}

¶ 18) 设 X 是这样的局部紧空间：在 X 中，每个相对紧开集的闭包仍是开集（Stone 空间；参见第二章 §1 习题 13）；设 \(\mu\) 是 X 上的正测度，其支集恰为 X，且使得 X 上每个有界可积数值函数都等价于一个有限连续函数（§4 习题 10）。

\begin{enumerate}
\def\labelenumi{\alph{enumi})}
\item
  证明在 X 中每个无处稠密集 N 都是局部可忽略的（对每个紧集 K，考察等价于 \(\varphi_{K\cap \overline{N}}\) 的连续函数）。
\item
  设 \(f\) 是 \(X\) 上的可测数值函数（有限或非有限），并设 \(g\) 是 \(X\) 上 \(\leqslant f\) 的最大下半连续函数。证明 \(f\) 与 \(g\) 局部几乎处处相等（注意，若 \(f\) 在紧集 \(K\) 上的限制连续，则 \(f\) 与 \(g\) 在 \(K\) 的内部相等，并利用 \(a\) ）。
\item
  由 b) 推出：在 X 中，每个关于 \(\mu\) 局部可忽略的集都是无处稠密集。特别地，X 中每个贫集都是无处稠密的。
\end{enumerate}

\begin{enumerate}
\def\labelenumi{\arabic{enumi})}
\setcounter{enumi}{18}
\item
  设 \(\mu\) 是局部紧空间 X 上的正测度。设 f 是 X 到完备度量空间 F 的 \(\mu\) -可测映射。为使 f 在 X 的紧子集 K 中能被可测阶梯函数一致逼近，必要且充分的是 \(f(\mathrm{K})\) 在 F 中相对紧。
\item
  设 \(X\) 是紧空间，\(\mu\) 是 \(X\) 上的正测度，\((f_n)\) 是一列 \(\mu\) -可测数值函数。证明下列性质等价：
\end{enumerate}

\begin{enumerate}
\def\labelenumi{(\roman{enumi})}
\item
  存在子序列 \((f_{n_k})\)（\((f_n)\) 的），它在 \(X\) 中几乎处处趋于 0；
\item
  存在一个由有限实数组成的序列 \((\lambda_{n})\)，使得
\end{enumerate}

\[
\limsup _ {n \to \infty} | \lambda_ {n} | > 0
\]

并且使得通项为 \(\lambda_{n}f_{n}(x)\) 的级数在 X 中几乎处处收敛；(iii) 存在一个由有限实数组成的序列 \((\lambda_n)\)，使得

\[
\sum_ {n = 1} ^ {\infty} | \lambda_ {n} | = + \infty
\]

并且使得通项为 \(\lambda_{n}f_{n}(x)\) 的级数在 X 中几乎处处绝对收敛。

（为证 (i) 蕴涵 (ii) 与 (iii)，利用 Egoroff 定理。为证 (iii) 蕴涵 (i)，证明 (iii) 蕴涵存在 X 的可测子集的递增序列 \((\mathbf{A}_k)\) 以及子序列 \((f_{n_k})\)（\((f_n)\) 的），使得 \(\mu (\mathbf{A}_k)\) 趋于 \(\mu (\mathbf{X})\) 且 \(\int |f_{n_k}\varphi_{\mathrm{A}_k}|d\mu\) 趋于 0。）

\begin{enumerate}
\def\labelenumi{\arabic{enumi})}
\setcounter{enumi}{20}
\item
  设 X 是局部紧空间，\(\mu\) 是 X 上的正测度，\((\mathrm{U}_{\alpha})_{\alpha\in\mathbf{A}}\) 是 X 的一个开覆盖。对每个 \(\alpha\in A\)，设 \(f_{\alpha}\) 是 \(U_{\alpha}\) 到某个集合 G 的映射。假设对每对指标 \(\alpha,\beta\)，点 \(x\in U_{\alpha}\cap U_{\beta}\) 中使 \(f_{\alpha}(x)\neq f_{\beta}(x)\) 者组成的集是局部 \(\mu\) -可忽略的。证明存在 X 到 G 的映射 f，使得对每个 \(\alpha\in A\)，\(x\in U_{\alpha}\) 中使 \(f(x)\neq f_{\alpha}(x)\) 者组成的集是局部 \(\mu\) -可忽略的。（先考虑 X 为紧的情形，用有限个集 \(\mathrm{U}_{\alpha}\) 覆盖 X；再借助第 9 号中的命题 14 过渡到一般情形。）
\item
  证明：若正测度 \(\mu\) 使得存在测度 \(>0\) 且可任意小的开集，则对每个 \(a > 0\)，由依测度收敛拓扑在集合 \(\mathbf{f} \in \mathcal{L}_{\mathrm{F}}^{p}\)（满足 \(\mathrm{N}_p(\mathbf{f}) \leqslant a\)）上诱导的拓扑严格粗于 \(p\) 阶平均收敛拓扑。
\item
  设 \((\mathbf{f}_{n})\) 是 \(S_{F}\) 中的一列函数，使得对每个可积集 A 以及每个子序列 \((\mathbf{f}_{n_{k}})\)（\((\mathbf{f}_{n})\) 的），存在 \((\mathbf{f}_{n_{k}})\) 的一个在 A 中几乎处处收敛于 0 的子序列；证明序列 \((\mathbf{f}_{n})\) 依测度收敛于 0。（用反证法。）
\end{enumerate}

¶ 24) 设 X 是 R 中的区间 {[}0,1{]}，μ 是 X 上的 Lebesgue 测度。证明对每个 Banach 空间 F，\(S_F\) 上的每个连续线性形式都恒等于零。（对 \(S_F\) 中 0 的每个邻域 V，证明存在整数 n \textgreater{} 0，使得 V + V + \ldots{} + V（n 项）包含一条直线。）

\begin{enumerate}
\def\labelenumi{\arabic{enumi})}
\setcounter{enumi}{24}
\tightlist
\item
  \begin{enumerate}
  \def\labelenumii{\alph{enumii})}
  \tightlist
  \item
    为使 \(S_{F}\) 的子集 H 是预紧的，必要且充分的是：对每个 \(\varepsilon > 0\) 以及每个可积集 \(A \subset X\)，存在紧集 \(M \subset F\) 以及 A 分成有限个可积集 \(A_{i}\) 的一个分割，使得对每个 \(f \in H\)，存在一个可积集 \(B \subset A\)（测度 \(\mu(B) \leqslant \varepsilon\)），具有下列性质：\(1^{\circ}\) \(f(A - B)\) 的每一点与 M 的距离 \(\leqslant \varepsilon\)；\(2^{\circ}\) 在每个集合 \(A_{i} \cap C B\) 中，f 的振幅 \(\leqslant \varepsilon\)。
  \end{enumerate}
\end{enumerate}

\begin{enumerate}
\def\labelenumi{\alph{enumi})}
\setcounter{enumi}{1}
\item
  证明：为使 \(\mathcal{L}_{\mathrm{F}}^{1}\) 的子集 H 相对拟紧，必要且充分的是它在 \(\mathcal{S}_{\mathrm{F}}\) 中预紧且等度可积。
\item
  证明：为使一列函数 \((\mathbf{f}_{n})\)（\(L_{F}^{1}\) 中的）是平均收敛拓扑下的 Cauchy 序列，必要且充分的是 \((\mathbf{f}_{n})\) 是依测度收敛拓扑下的 Cauchy 序列，且诸 \(f_{n}\) 组成的集等度可积。
\end{enumerate}

\(d\) ) 把这些性质推广到空间 \(\mathcal{L}_{\mathrm{F}}^{p}\)（\(p > 1\)）

\begin{enumerate}
\def\labelenumi{\arabic{enumi})}
\setcounter{enumi}{25}
\tightlist
\item
  设 \(f\) 是定义在 \(\mathbf{R}\) 上的数值函数。证明若对每个 \(x \in \mathbf{R}\) 有 \(\liminf_{y \to x, y \geqslant x} f(y) \geqslant f(x)\)，则 \(f\) 对每个测度 \(\mu\) 都是 \(\mu\) -可测的（\(\mathbf{R}\) 上的）。
\end{enumerate}

¶ 27) 设 X 是局部紧空间，\(\mu\) 是 X 上的实测度，K 是 X 的紧子集，\((f_n)\) 是定义在 K 上且分隔 K 的点的有界 \(\mu\) -可测数值函数序列。证明：对定义在 K 上的每个有界 \(\mu\) -可测数值函数 g 以及每个 \(\varepsilon > 0\)，存在一个多项式 \(h = \mathrm{P}((f_n))\)（\(f_n\) 的）以及 K 的紧子集 \(\mathrm{K}_1\)，使得 \(\|h\| \leqslant 2\|g\|\)、\(|\mu|(\mathrm{K} - \mathrm{K}_1) \leqslant \varepsilon\)，且对每个 \(x \in \mathrm{K}_1\)，\(|h(x) - g(x)| \leqslant \varepsilon\)。（考察映射 \(x \mapsto (f_n(x))\)（K 到 \(\mathbf{R}^{\mathbf{N}}\) 中的）及其像的闭包，并适当应用 Stone--Weierstrass 定理。）由此推出，对 \(1 \leqslant p < +\infty\)，\(f_n\) 的多项式所成的集在 \(\mathcal{L}^p(\mathrm{K})\) 中稠密。当把序列 \((f_n)\) 换成分隔 K 的点的不可数有界可测函数族时，这些性质还成立吗？

\begin{enumerate}
\def\labelenumi{\arabic{enumi})}
\setcounter{enumi}{27}
\tightlist
\item
  设 X 是局部紧空间，\(\mu\) 是 X 上的正测度，P 是形如 \(f \geqslant 0\)（有限或非有限）的数值函数（定义在 X 上且 \(\mu\) -可测）所成的集。设 \(\lambda\) 是定义在 P 上、取值 \(\geqslant 0\)（有限或非有限）的函数，它正齐次、递增且凸（参见第一章第 1 号），此外还满足：\(1^{\circ} \lambda(f) = 0\) 对每个 \(\mu\) -可忽略函数 \(f \geqslant 0\) 成立；\(2^{\circ}\) 对 P 中函数的每个递增序列 \((f_{n})\)，有 \(\lambda(\sup f_{n}) = \sup_{n} \lambda(f_{n})\)。对每个 Banach 空间 F，用 \(L_{F}^{\lambda}\) 表示 X 到 F 的 \(\mu\) -可测映射 f（使 \(\lambda(|\mathbf{f}|)\) 有限者）所成的集。证明 \(\lambda(|\mathbf{f}|)\) 是 \(L_{F}^{\lambda}\) 上的一个半范数，且 \(L_{F}^{\lambda}\) 对由此半范数定义的拓扑是完备的。
\end{enumerate}

¶ 29) 设 \(\mu\) 是局部紧空间 X 上的测度。对每个定义在 X 上的 \(\mu\) -可测数值函数 \(f \geqslant 0\)，映射 \(f_{\mu}' : t \mapsto |\mu|^*(f([t, +\infty]))\)（\(\mathbf{R}_+\) 到 \(\overline{\mathbf{R}}_{+}\) 的）是递减且右连续的。若 \(\nu\) 是第二个局部紧空间 Y 上的测度，且 \(g\) 是定义在 Y 上的 \(\nu\) -可测数值函数 \(\geqslant 0\)，则称 \(f\) 与 \(g\) 等可测（分别关于 \(\mu\) 与 \(\nu\)），如果 \(f_{\mu}^{\prime} = g_{\nu}^{\prime}\)。用 \(f^{*}\)（或 \(f_{\mu}^{*}\)）表示定义在 \(\mathbf{R}_{+}\) 上的函数，它对每个 \(s \in \mathbf{R}_{+}\)，在 \(\overline{\mathbf{R}}_{+}\) 中等于由数 \(a\)（满足 \(f_{\mu}^{\prime}(a) \geqslant s\) 者）组成之集的上确界（若该集为空，则上确界等于 0）。

\begin{enumerate}
\def\labelenumi{\alph{enumi})}
\item
  证明函数 \(f^{*}\) 是递减的、在 \(\mathbf{R}_{+}\) 中左连续，且与 \(f\) 等可测（关于测度 \(\mu\) 与 \(\mathbf{R}_{+}\) 上的 Lebesgue 测度）；\(f^{*}\) 称为 \(f\) 的递减重排。
\item
  若 \(0 \leqslant f \leqslant g\) 是两个 \(\mu\) -可测函数（\(X\) 上的），证明 \(f^* \leqslant g^*\)。若 \((f_n)\) 是由 \(\mu\) -可测函数 \(\geqslant 0\)（\(X\) 上的）组成的递增序列，且 \(f = \sup_{n} f_n\)，证明
\end{enumerate}

即 \(f^{*}(s) = \sup_{n}f_{n}^{*}(s)\) 在 \(f^{*}\) 连续的每一点处成立。

\begin{enumerate}
\def\labelenumi{\alph{enumi})}
\setcounter{enumi}{2}
\item
  证明对每个 \(\mu\) -可测数值函数 \(f \geqslant 0\)，有 \(|\mu|^{*}(f) = \int^{*} f^{*}(s) ds\)。（先考虑阶梯函数的情形，然后利用 \(b\) ）。）
\item
  设 w 是定义在 \(R_{+}\) 上的数值函数，它是有限的（除点 s = 0 外）且递减。对每个定义在 X 上的 \(\mu\) -可测数值函数 \(f \geqslant 0\)，令
\end{enumerate}

\[
\lambda (f) = \left(\int^ {*} w (s) \bigl (f ^ {*} (s) \bigr) ^ {p} d s\right) ^ {1 / p} \quad (1 \leqslant p <   + \infty).
\]

证明此函数满足习题 28 的条件。（为证其凸性，先利用 c) 考察 w 为递减阶梯函数的情形，然后取极限处理一般情形。）

\begin{enumerate}
\def\labelenumi{\alph{enumi})}
\setcounter{enumi}{4}
\tightlist
\item
  沿用 \(d\) ) 的记号，用 \(\mathcal{L}_{\mathrm{F}}^{p,w}\) 代替 \(\mathcal{L}_{\mathrm{F}}^{\lambda}\) 来记，并令 \(\mathrm{N}_{p,w}(\mathbf{f}) = \lambda (|\mathbf{f}|)\)（对 X 到 F 的每个 \(\mu\) -可测映射 \(\mathbf{f}\)）。证明：若 \(\mathbf{f} \in \mathcal{L}_{\mathrm{F}}^{p,w}\)，且对每个 \(n\)，用 \(\mathbf{f}_n\) 表示等于 \(\mathbf{f}\)（若 \(|\mathbf{f}| \leqslant n\)）而等于 \(nf / |\mathbf{f}|\)（若 \(|\mathbf{f}| > n\)）的函数，则序列 \((|\mathbf{f}_n|^*)\) 几乎处处趋于 \(|\mathbf{f}|^*\)，且有 \(\lim_{n \to \infty} \mathrm{N}_{p,w}(\mathbf{f} - \mathbf{f}_n) = 0\)。若 \(\mathbf{f}\) 是 \(\mu\) -可忽略的，则 \(\mathrm{N}_{p,w}(\mathbf{f}) = 0\)。
\end{enumerate}

¶ 30) 设 \(\Phi(t, u, v)\) 是在 \(0 \leqslant t \leqslant 1\)、\(u \geqslant 0\)、\(v \geqslant 0\) 上有限且连续的数值函数。对每个数值函数 \(f \geqslant 0\)（定义在 I = {[}0,1{]} 上、关于 Lebesgue 测度可测），用 \(f^*\) 表示 \(f\) 的递减重排（习题 29）。为使对每对函数 \(f \geqslant 0\)、\(g \geqslant 0\)（在 I 上（关于 Lebesgue 测度）可测且有界），都有

\[
\int_ {0} ^ {1} \Phi (t, f (t), g (t)) d t \leqslant \int_ {0} ^ {1} \Phi (t, f ^ {*} (t), g ^ {*} (t)) d t,\tag{1}
\]

必要且充分的是函数 \(\Phi\) 满足以下三个条件：

\[
\Phi (t, u + h, v + h) - \Phi (t, u + h, v) - \Phi (t, u, v + h) + \Phi (t, u, v) \geqslant 0\tag{2}
\]

对一切 \(t \in I, u \geqslant 0, v \geqslant 0\) 以及 \(h \geqslant 0\) ；

\[
\begin{array}{r l} \int_ {0} ^ {\delta} & (\Phi (a + \delta + t, u, v) - \Phi (a + \delta + t, u + h, v) \\ & + \Phi (a + t, u + h, v) - \Phi (a + t, u, v)) d t \geqslant 0 \end{array}\tag{3}
\]

以及

\[
\begin{array}{l} \int_ {0} ^ {\delta} \big (\Phi (a + \delta + t, u, v) - \Phi (a + \delta + t, u, v + h) \\ \qquad \qquad \qquad + \Phi (a + t, u, v + h) - \Phi (a + t, u, v) \big) d t \geqslant 0 \end{array}\tag{4}
\]

对一切 \(a, h, \delta\)（满足 \(0 \leqslant a \leqslant 1 - 2\delta\)）以及 \(h \geqslant 0\)。

（为证 (3) 与 (4) 是必要的，取适当的阶梯函数作为 \(f\) 与 \(g\)；再由 (3) 与 (4) 推出 (2)。为证这些条件是充分的，首先利用 \(\Phi\) 的连续性由 (2) 推出

\[
\Phi (t, u + h, v + k) - \Phi (t, u + h, v) - \Phi (t, u, v + k) + \Phi (t, u, v) \geqslant 0\tag{5}
\]

对一切 \(t \in \mathrm{I}\)、\(u \geqslant 0\)、\(v \geqslant 0\)、\(h \geqslant 0\)、\(k \geqslant 0\)。然后利用 (3) 与 (4)，对 \(f\) 与 \(g\) 为间断点（在 \(\mathrm{I}\) 中）形如 \(r / n\)（\(0 \leqslant r \leqslant n\)）的阶梯函数证明 (1)，方法是依次比较下列积分

\[
\int_ {0} ^ {1} \Phi (t, f _ {i} (t), g _ {i} (t)) d t
\]

以及

\[
\int_ {0} ^ {1} \Phi (t, f _ {j} (t), g _ {j} (t)) d t,
\]

其中 \(f_{i}\) 与 \(g_{i}\) 同 \(f_{j}\) 与 \(g_{j}\) 只在两个相邻区间 \([(r - 1) / n, r / n]\) 与 \([r / n, (r + 1) / n]\) 上不同。最后，以适当的方式取极限，以便在一般情形证明 (1)。）

若 \(\Phi\) 二次连续可微，则条件 (2)、(3) 与 (4) 分别等价于

\[
\frac {\partial^ {2} \Phi}{\partial u \partial v} \geqslant 0, \quad \frac {\partial^ {2} \Phi}{\partial t \partial u} \leqslant 0, \quad \frac {\partial^ {2} \Phi}{\partial t \partial v} \leqslant 0.
\]

把上述结果推广到任意多个变量的函数 \(\Phi(t, u_1, \ldots, u_m)\)。

\setcounter{exercisegroup}{6}
\edef\CurrentExerciseGroup{\arabic{exercisegroup}}
